\begin{exo}{}
	Le tableau ci-après donne les taux d'évolution successifs du nombre d'objets connectés dans le monde entre les années $2011$ et $\np{2015}$.
	La cellule \textbf{B6} est au format pourcentage écrit avec une décimale (\textit{source : IDATE DigiWorld, Institut de l'audiovisuel et des télécoms en Europe}).


	\begin{center}
		\def\arraystretch{1.5}
		\begin{tabular}{|c|c|c|c|c|c|}
			\hline
			                     & \cellcolor{gray!30}A & \cellcolor{gray!30}B  & \cellcolor{gray!30}C & \cellcolor{gray!30}D & \cellcolor{gray!30}E \\ \hline
			\cellcolor{gray!30}1 & \textbf{Année}       & 2012                  & 2013                 & 2014                 & 2015                 \\ \hline
			\cellcolor{gray!30}2 & \textbf{Taux}        & $57,89$\%             & $46,67$\%            & $40,91$\%            & $35,48$\%            \\ \hline
			\cellcolor{gray!30}3 & \textbf{C.M.}        &                       &                      &                      &                      \\ \hline
			\cellcolor{gray!30}4 &                      &                       &                      &                      &                      \\ \hline
			\cellcolor{gray!30}5 & $CM_g$               & \multicolumn{4}{l|}{}                                                                      \\ \hline
			\cellcolor{gray!30}6 & $T_g$                & \multicolumn{4}{l|}{}                                                                      \\ \hline
		\end{tabular}
	\end{center}


	\begin{enumerate}
		\item Quelle formule faut-il entrer en \textbf{B3} puis recopier vers la droite pour calculer le coefficient multiplicateur associé aux taux d'évolution du dessus ?
		\item Quelles formules faut-il entrer en \textbf{B5} et \textbf{B6} pour calculer le coefficient multiplicateur global et le taux d'évolution global ?
		\item Recopier et compléter le tableau.
	\end{enumerate}
\end{exo}
